\chapter{Vector and Tensor Perturbations on FLRW}

The scalar--vector--tensor split separates decaying vector modes from propagating tensor modes on FLRW.  The discrete tensor sector preserves the continuum causal structure in the refinement limit.

\allowdisplaybreaks

\section{Scalar--vector--tensor split}
In conformal time,
\begin{equation}
\dd s^2=a(\eta)^2
\left[
-\dd\eta^2
+2S_i\dd\eta\dd x^i
+\left(\delta_{ij}+2\partial_{(i}F_{j)}+h_{ij}\right)
\dd x^i\dd x^j
\right],
\end{equation}
with
\begin{align}
\partial^iS_i&=0,\\
\partial^iF_i&=0,\\
\partial^ih_{ij}&=0,\\
h^i{}_i&=0.
\end{align}
Under a transverse spatial gauge shift \(L_i\),
\begin{align}
S_i&\mapsto S_i-L_i',\\
F_i&\mapsto F_i-L_i.
\end{align}
Therefore
\begin{equation}
\boxed{
\Sigma_i:=S_i-F_i'
}
\end{equation}
is gauge invariant, because
\begin{align}
\Sigma_i
&\mapsto
(S_i-L_i')-(F_i-L_i)'\\
&=S_i-L_i'-F_i'+L_i'\\
&=\Sigma_i.
\end{align}

\section{Source-free vector sector}
In the absence of vector anisotropic stress,
\begin{equation}
\boxed{
(a^2\Sigma_i)'=0.
}
\end{equation}
Integrating once gives
\begin{align}
(a^2\Sigma_i)'&=0,\\
a^2\Sigma_i&=C_i,\\
\Sigma_i(\eta)&=\frac{C_i}{a(\eta)^2}.
\end{align}
For a matter background \(a\propto\eta^2\),
\begin{equation}
\Sigma_i\propto\eta^{-4}.
\end{equation}
A midpoint update of
\begin{equation}
\Sigma'=-\frac{4}{\eta}\Sigma
\end{equation}
converges with
\begin{equation}
p_{\mathrm{vector}}=2.052951.
\end{equation}
The fully matter-sourced vector constraint and vorticity hierarchy are \text{\normalfont not derived here}.

\section{Tensor action and equation}
The quadratic tensor action is
\begin{equation}
\boxed{
S_T^{(2)}
=\frac{M_{\mathrm{Pl}}^2}{8}
\int a^2
\left[
(h_{ij}')^2-(\partial_kh_{ij})^2
\right]
\dd\eta\dd^3x.
}
\end{equation}
For a transverse-traceless variation,
\begin{align}
\delta S_T^{(2)}
&=\frac{M_{\mathrm{Pl}}^2}{4}
\int a^2
\left[
h_{ij}'\delta h_{ij}'
-\partial_kh_{ij}\partial_k\delta h_{ij}
\right]\dd\eta\dd^3x\\
&=-\frac{M_{\mathrm{Pl}}^2}{4}
\int
\left[
(a^2h_{ij}')'-a^2\nabla^2h_{ij}
\right]
\delta h_{ij}\,\dd\eta\dd^3x.
\end{align}
Since
\begin{equation}
(a^2h_{ij}')'
=a^2h_{ij}''+2aa'h_{ij}'
=a^2(h_{ij}''+2\mathcal Hh_{ij}'),
\end{equation}
the vacuum Euler--Lagrange equation is
\begin{equation}
\boxed{
h_{ij}''+2\mathcal Hh_{ij}'-\nabla^2h_{ij}=0,
\qquad
\mathcal H:=a'/a.
}
\end{equation}
With transverse-traceless anisotropic stress,
\begin{equation}
h_{ij}''+2\mathcal Hh_{ij}'-\nabla^2h_{ij}
=\frac{2a^2}{M_{\mathrm{Pl}}^2}\Pi_{ij}^{\mathrm{TT}}.
\end{equation}
For each Fourier polarization, set
\begin{equation}
\mu_\lambda:=\frac{aM_{\mathrm{Pl}}}{2}h_\lambda,
\qquad
h_\lambda=\frac{2\mu_\lambda}{aM_{\mathrm{Pl}}}.
\end{equation}
Then
\begin{align}
h_\lambda'
&=\frac{2}{aM_{\mathrm{Pl}}}
(\mu_\lambda'-\mathcal H\mu_\lambda),\\
h_\lambda''
&=\frac{2}{aM_{\mathrm{Pl}}}
\left[
\mu_\lambda''-2\mathcal H\mu_\lambda'
+(\mathcal H^2-\mathcal H')\mu_\lambda
\right].
\end{align}
Substitution into
\(h_\lambda''+2\mathcal Hh_\lambda'+k^2h_\lambda=0\)
cancels the first-derivative terms and gives
\begin{align}
0
&=\mu_\lambda''+
\left[k^2-(\mathcal H'+\mathcal H^2)\right]\mu_\lambda\\
&=\mu_\lambda''+
\left(k^2-\frac{a''}{a}\right)\mu_\lambda.
\end{align}
Thus
\begin{equation}
\boxed{
\mu_\lambda''+
\left(k^2-\frac{a''}{a}\right)\mu_\lambda=0.
}
\end{equation}

\section{Projector-lattice tensor sector}
For cubic spacing \(h\), the centered projector Laplacian has symbol
\begin{align}
\widehat k_h^2
&=\frac4{h^2}\sum_{r=1}^3
\sin^2\left(\frac{k_rh}{2}\right)\\
&=\sum_{r=1}^3
\left[
k_r^2-\frac{h^2k_r^4}{12}
+\frac{h^4k_r^6}{360}
+\order(h^6)
\right]\\
&=k^2-\frac{h^2}{12}\sum_{r=1}^3k_r^4
+\frac{h^4}{360}\sum_{r=1}^3k_r^6
+\order(h^6).
\end{align}
Accordingly,
\begin{align}
\mu_{\lambda,h}''+
\left(\widehat k_h^2-\frac{a''}{a}\right)
\mu_{\lambda,h}&=0.
\end{align}
Radiation era:
\begin{align}
a(\eta)&\propto\eta,
&a''&=0,\\
\mu_\kappa(\eta)&=\frac{\ee^{-\ii\kappa\eta}}{\sqrt{2\kappa}},
&h_\kappa&=\frac{2\mu_\kappa}{aM_{\mathrm{Pl}}}.
\end{align}
Thus
\begin{equation}
a|h_\kappa|=\mathrm{constant}.
\end{equation}
Matter era:
\begin{align}
a(\eta)&\propto\eta^2,\\
h_\kappa(\eta)
&=3\frac{j_1(\kappa\eta)}{\kappa\eta},\\
h_\kappa''+\frac4\eta h_\kappa'+\kappa^2h_\kappa&=0.
\end{align}
The numerical orders are
\begin{align}
p_T^{\mathrm{radiation}}&=1.999618,\\
p_{T,\mathrm R}^{\mathrm{radiation}}&=3.998798,\\
p_T^{\mathrm{matter}}&=2.002687,\\
p_{T,\mathrm R}^{\mathrm{matter}}&=3.996376.
\end{align}
The canonical Wronskian and amplitude residuals satisfy
\begin{align}
\varepsilon_{\mathrm W}&=4.440892\times10^{-16},\\
\varepsilon_{a|h|}&=4.440892\times10^{-16}.
\end{align}

\section{Propagation distance}
For geometric-optics tensor propagation, the canonical variable obeys the locally flat radial wave law
\begin{equation}
\mu\sim\frac{\ee^{-\ii k\eta}}{r}.
\end{equation}
Using \(h=2\mu/(aM_{\mathrm{Pl}})\) then gives
\begin{equation}
h\sim\frac{1}{ar},
\end{equation}
up to the source normalization and polarization amplitude.
After source-frame normalization,
\begin{equation}
D_L^{\mathrm{GW}}=D_L^{\mathrm{EM}}
\end{equation}
in the unmodified propagation sector. Source generation, binary calibration, and modified-friction alternatives are \text{\normalfont not derived here}.

\section{Open extensions}
\begin{nextsteps}
\item Add the matter-sourced vector constraint and derive the complete vorticity hierarchy.
\item Include cosmological anisotropic stress and its contribution to vector and tensor propagation.
\item Derive the primordial tensor normalization from the selected quantum initial state.
\item Compute the leading nonlinear scalar--vector--tensor couplings.
\item Incorporate source generation, binary calibration, and standard-siren observables in the tensor sector.
\end{nextsteps}
